\documentclass[../../../include/open-logic-section]{subfiles}

\begin{document}

\olfileid{sth}{ordinals}{basic}
\olsection{ترتیبی عدداں دیاں بنیادی خاصیتاں}

اسیں ویکھیا سی کہ پہلے کجھ ترتیبی عدد قدرتی عدد نیں۔ ترتیبی عدداں دا
نظریہ بناؤن دی وڈی وجہ اوہ استقرائی اصول ودھاؤنا اے جیہڑا قدرتی عدداں
لئی قائم اے۔ اسیں کجھ بنیادی نتیجیاں دی لڑی راہیں اوس اصول تک پہنچاں گے۔

\begin{lem}\ollabel{ordmemberord}
ترتیبی عدد دا ہر !!{element} آپ اک ترتیبی عدد اے۔
\end{lem}

\begin{proof}
منّ لو $\alpha$ ترتیبی عدد اے تے $b \in \alpha$۔ چونکہ $\alpha$
متعدی اے، $b \subseteq \alpha$۔ سو $\in$، $b$ نوں حُسنِ ترتیب وچ
رکھدا اے، کیوں جے $\in$، $\alpha$ نوں حُسنِ ترتیب وچ رکھدا اے۔

ایہ ویکھن لئی کہ $b$ متعدی اے، منّ لو $x \in c \in b$۔ سو $c \in \alpha$،
کیوں جے $b \subseteq \alpha$۔ فیر چونکہ $\alpha$ متعدی اے،
$c \subseteq \alpha$، تے ایس لئی $x \in \alpha$۔ سو $x, c, b \in \alpha$۔
پر $\in$، $\alpha$ نوں حُسنِ ترتیب وچ رکھدا اے، ایس لئی $\in$، $\alpha$
اُتے متعدی تعلق اے، \olref[sth][ordinals][wo]{wo:strictorder} نال۔
سو چونکہ $x \in c \in b$، ایس لئی $x \in b$۔ عام صورت وچ $c \subseteq b$۔
\end{proof}

\begin{cor}\ollabel{ordissetofsmallerord}
$\alpha = \Setabs{\beta \in \alpha}{\beta \text{ ترتیبی عدد اے}}$،
ہر ترتیبی عدد~$\alpha$ لئی۔
\end{cor}

\begin{proof}
\olref{ordmemberord} توں سدھا نکلدا اے۔
\end{proof}

اگلے دو اہم نتیجیاں، \olref{ordinductionschema} تے \olref{ordtrichotomy}،
دی موٹی گل ایہ اے کہ رکنیت آپ ترتیبی عدداں نوں حُسنِ ترتیب وچ رکھدی اے:

\begin{thm}[لامتناہی درجیاں اُتے استقرا]\ollabel{ordinductionschema}
ہر فارمولے $\phi(x)$ لئی:
\[
	\text{ جے }\exists \alpha \phi(\alpha)\text{، تاں }\exists \alpha(\phi(\alpha)
	\land  (\forall \beta \in \alpha) \lnot \phi(\beta))
\]
ایتھے دِتے مقداری نشان ان کہے طور اُتے ترتیبی عدداں تک محدود نیں۔
\end{thm}

\begin{proof}
%\olref{ordinductionschema1}. 
منّ لو $\phi(\alpha)$ کسے ترتیبی عدد $\alpha$ لئی۔ جے $ (\forall \beta
\in \alpha) \lnot \phi(\beta)$، تاں ثبوت پورا ہویا۔ نہیں تاں چونکہ
$\alpha$ ترتیبی عدد اے، اوہدے وچ $\in$ دے لحاظ نال اک سبھ توں چھوٹا
!!{element} اے جیہڑا $\phi$ پورا کردا اے، تے \olref{ordmemberord} نال
ایہ ترتیبی عدد اے۔
%	\olref{ordinductionschema2}. Suppose there is some ordinal
%	$\gamma$ such that $\lnot\phi(\gamma)$. Then by
%	\olref{ordinductionschema1}, there is an $\in$-minimal ordinal
%	$\alpha$ for which $\lnot\phi(\alpha)$. So $(\forall \beta <
%	\alpha) \phi(\beta)$, rendering the antecedent of the conditional
%	false.
\end{proof}
\noindent
دھیان دیو کہ \olref{ordinductionschema} نوں ایس خاکے دے روپ وچ وی لکھ سکدے آں:
\[
\text{ جے }\forall \alpha((\forall \beta \in \alpha)\phi(\beta) \lif 
\phi(\alpha))\text{، تاں }\forall \alpha\phi(\alpha)
\]
بس $\lnot\phi(\alpha)$ نوں \olref{ordinductionschema} وچ لینا اے تے فیر بنیادی منطقی تبدیلیاں کرنیاں نیں۔

\begin{thm}[تِن صورتاں دا اصول]\ollabel{ordtrichotomy}
$\alpha \in \beta \lor \alpha = \beta \lor \beta \in \alpha$، کسے وی
دو ترتیبی عدداں $\alpha$ تے $\beta$ لئی۔
\end{thm}

\begin{proof}
ثبوت دوہرے استقرا نال اے، یعنی \olref{ordinductionschema} نوں دو واری
ورت کے۔ آکھو کہ $x$، $y$ نال \emph{پرکھیا جا سکدا اے} اُدوں تے صرف اُدوں
جدوں $x \in y \lor x = y \lor y \in x$۔

استقرا لئی منّ لو کہ~$\alpha$ دے ہر ترتیبی عدد دی \emph{ہر} ترتیبی
عدد نال پرکھ ہو سکدی اے۔ ہور استقرا لئی منّ لو کہ $\alpha$ دی~$\beta$
دے ہر ترتیبی عدد نال پرکھ ہو سکدی اے۔ اسیں وکھاواں گے کہ $\alpha$
دی~$\beta$ نال پرکھ ہو سکدی اے۔~$\beta$ اُتے استقرا نال نکلے گا کہ $\alpha$
دی ہر ترتیبی عدد نال پرکھ ہو سکدی اے۔ فیر~$\alpha$ اُتے استقرا نال،
\emph{ہر} ترتیبی عدد دی \emph{ہر} ترتیبی عدد نال پرکھ ہو سکدی اے،
جیہڑی سانوں چاہیدی سی۔ ایہ کافی اے کہ $\alpha \notin \beta$ تے
$\beta \notin \alpha$ منّ کے وکھائیے کہ $\alpha = \beta$۔

ایہ وکھاؤن لئی کہ $\alpha \subseteq \beta$، $\gamma \in \alpha$ مقرر کرو۔
\olref{ordmemberord} نال ایہ اک ترتیبی عدد اے۔ سو پہلے استقرائی مفروضے
نال $\gamma$ دی $\beta$ نال پرکھ ہو سکدی اے۔ پر جے یا تاں $\gamma = \beta$
یا $\beta \in \gamma$، تاں $\beta \in \alpha$ (لوڑ پئے تاں ایس گل نوں
ورت کے کہ $\alpha$ متعدی اے)۔ ایہ ساڈے مفروضے دے اُلٹ اے۔ سو
$\gamma \in \beta$۔ عام صورت وچ $\alpha \subseteq \beta$۔

دوجے استقرائی مفروضے نال عین ایہو جیہی دلیل وکھاندی اے کہ
$\beta \subseteq \alpha$۔ سو $\alpha = \beta$۔
\end{proof}\noindent ایس لئی کدی اسیں $\alpha <\beta$ لکھاں گے،
$\alpha \in \beta$ دی تھاں، کیوں جے $\in$ ترتیب دے تعلق وانگ کم کردا اے۔
ایہدی کوئی ڈونگھی وجہ نہیں، بس ایہ جانیا پچھانیا اے تے
$\alpha \leq \beta$ لکھنا $\alpha \in \beta \lor \alpha = \beta$
لکھن نالوں سوکھا اے۔\footnote{اسیں $\alpha \mathrel{\underline{\in}} \beta$
لکھ سکدے ساں، پر ایہ بالکل غیر معیاری ہوندا۔}

ساڈے پچھلے نتیجیاں دے دو سدھے حاصل ایہ نیں۔ پہلا ساڈے نویں نشاناں نوں کم وچ لیاندا اے:

\begin{cor}\ollabel{ordordered}
جے $\exists \alpha\phi(\alpha)$، تاں $\exists \alpha(\phi(\alpha)
\land \forall \beta(\phi(\beta) \lif \alpha \leq \beta))$۔
نالے کسے وی ترتیبی عدداں $\alpha, \beta, \gamma$ لئی،
$\alpha \notin \alpha$ تے $\alpha \in \beta \in \gamma \lif \alpha \in \gamma$
دونیں قائم نیں۔
\end{cor}

\begin{proof}
بالکل \olref[wo]{wo:strictorder} وانگ۔
\end{proof}

\begin{prob}
\olref[sth][ordinals][basic]{ordtrichotomy} دے ثبوت وچ ’’عین ایہو جیہی دلیل‘‘
نوں پورا کرو۔
\end{prob}

\begin{cor}\ollabel{corordtransitiveord}
$A$ ترتیبی عدد اُدوں تے صرف اُدوں اے جدوں $A$ ترتیبی عدداں دا متعدی سیٹ اے۔
\end{cor}

\begin{proof}
\emph{کھبے توں سجے پاسے۔} \olref{ordmemberord} نال۔
\emph{سجے توں کھبے پاسے۔} جے $A$ ترتیبی عدداں دا متعدی سیٹ اے، تاں
$\in$، $A$ نوں حُسنِ ترتیب وچ رکھدا اے،
\olref{ordinductionschema} تے \olref{ordtrichotomy} نال۔
\end{proof}

اسیں \olref{ordinductionschema} تے \olref{ordtrichotomy} دا مطلب ایہ دسیا سی
کہ $\in$ ترتیبی عدداں نوں حُسنِ ترتیب وچ رکھدا اے۔ پر ہِیٹھلے نتیجے کر کے
ایہو جیہے دعوے بارے سانوں \emph{بڑی احتیاط} کرنی پئے گی:

\begin{thm}[بورالی فورتی دا پیراڈوکس]\ollabel{buraliforti}
سارے ترتیبی عدداں دا کوئی سیٹ نہیں۔
\end{thm}

\begin{proof}
تضاد کڈھن لئی منّ لو $O$ سارے ترتیبی عدداں دا سیٹ اے۔ جے
$\alpha \in \beta \in O$، تاں $\alpha$ ترتیبی عدد اے، \olref{ordmemberord}
نال، سو $\alpha \in O$۔ ایس لئی $O$ متعدی اے، تے \olref{corordtransitiveord}
نال $O$ ترتیبی عدد اے۔ ایس لئی $O \in O$، جیہڑا
\olref{ordordered} دے اُلٹ اے۔
\end{proof}
\noindent
ایہ نتیجہ \citeauthor{Burali-Forti1897} دے ناں اُتے اے۔ پر 1899 وچ ڈیڈیکنڈ
نوں اک خط وچ کانتور نے سبھ توں پہلاں صاف ویکھیا کہ سارے ترتیبی عدداں دا
سیٹ منّن وچ \emph{تضاد} اے۔ جیویں فان ہائینورٹ دسدا اے:
\begin{quote}
  بورالی فورتی آپ ایس تضاد نوں \emph{برہانِ خلف} نال ایس نتیجے دا ثبوت
  سمجھدا سی کہ ترتیبی عدداں دی قدرتی ترتیب بس اک جزوی ترتیب اے۔
  \citep[p.~105]{Heijenoort1967}
\end{quote}
بورالی فورتی دی غلطی اک پاسے رکھ کے، پچھلی گل ایویں سمیٹ سکدے آں۔
ترتیبی عدد سیٹ نیں۔ رکنیت ہر اک نوں حُسنِ ترتیب وچ رکھدی اے تے
ساریاں نوں رلا کے وی حُسنِ ترتیب وچ رکھدی اے (بھانویں سارے رل کے سیٹ
نہیں بناندے)۔

گل پوری کرن لئی ترتیبی عدداں دیاں کجھ ہور بنیادی خاصیتاں ہِیٹھاں نیں،
جیہڑیاں \olref{ordinductionschema} تے \olref{ordtrichotomy} توں نکلدیاں نیں۔

\begin{prop}
ترتیبی عدداں دی ہر سختی نال گھٹدی لڑی محدود اے۔
\end{prop}

\begin{proof}
ترتیبی عدداں دی کسے وی لامتناہی سختی نال گھٹدی لڑی
$\alpha_0 > \alpha_1 > \alpha_2 > \ldots$ وچ $<$ دے لحاظ نال
کوئی ایسا رکن نہیں جیہدے توں ہور کوئی رکن چھوٹا نہ ہووے۔
ایہ \olref{ordinductionschema} دے اُلٹ اے۔
\end{proof}

\begin{prop}\ollabel{ordinalsaresubsets}
$\alpha \subseteq \beta \lor \beta \subseteq \alpha$، کسے وی
ترتیبی عدداں $\alpha, \beta$ لئی۔
\end{prop}

\begin{proof}
جے $\alpha \in \beta$، تاں $\alpha \subseteq \beta$، کیوں جے $\beta$
متعدی اے۔ ایسے طرح جے $\beta \in \alpha$، تاں $\beta \subseteq \alpha$۔
تے جے $\alpha = \beta$، تاں $\alpha \subseteq \beta$ تے $\beta \subseteq \alpha$۔
سو \olref{ordtrichotomy} نال ثبوت پورا ہویا۔
\end{proof}

\begin{prop}\ollabel{ordisoidentity}
$\alpha = \beta$ اُدوں تے صرف اُدوں اے جدوں $\ordeq{\alpha}{\beta}$،
کسے وی ترتیبی عدداں $\alpha, \beta$ لئی۔
\end{prop}

\begin{proof}
ترتیبی عدد حُسنِ ترتیب والیاں ساختاں نیں۔ سو ایہ تِن صورتاں دے اصول
(\olref[basic]{ordtrichotomy}) تے \olref[iso]{wellordnotinitial}
توں سدھا نکلدا اے۔
\end{proof}

\begin{prob}\ollabel{probunionordinalsordinal}
ثابت کرو کہ جے $X$ دا ہر رکن ترتیبی عدد اے، تاں $\bigcup X$ وی ترتیبی عدد اے۔
\end{prob}

\end{document}
